% Bachelor Thesis Exposé — FINAL (revised per supervisor feedback, July 2026)
% Changes from v1: (1) title now includes hardware name "LIMO Cobot";
%                  (2) Gantt redrawn for 3-month period Jul 1 – Oct 1, 2026
%                      with more realistic experiment allocation.
% Compile: pdflatex expose_final.tex

\documentclass[11pt,a4paper]{article}

\usepackage[utf8]{inputenc}
\usepackage[T1]{fontenc}
\usepackage[margin=2.5cm]{geometry}
\usepackage{setspace}
\usepackage{enumitem}
\usepackage{xcolor}
\usepackage{tabularx}
\usepackage{hyperref}
\usepackage{fancyhdr}
\usepackage{titlesec}
\usepackage{booktabs}
\usepackage{tcolorbox}
\usepackage{amssymb}
\usepackage{pgfgantt}
\usepackage{array}
\usepackage{colortbl}
\usepackage{multirow}

% ── Colours ──────────────────────────────────────────────────────────────────
\definecolor{primary}{RGB}{37,99,235}
\definecolor{heading}{RGB}{30,41,59}
\definecolor{body}{RGB}{51,65,85}
\definecolor{accent}{RGB}{220,38,38}
\definecolor{boxbg}{RGB}{248,250,252}
\definecolor{border}{RGB}{226,232,240}
\definecolor{ganttblue}{RGB}{37,99,235}
\definecolor{ganttgrey}{RGB}{100,116,139}

% ── Page style ────────────────────────────────────────────────────────────────
\pagestyle{fancy}
\fancyhf{}
\fancyfoot[C]{\thepage}
\renewcommand{\headrulewidth}{0pt}

% ── Section formatting ────────────────────────────────────────────────────────
\titleformat{\section}{\Large\bfseries\color{heading}}{}{0em}{}
    [\vspace{-0.3em}\rule{\textwidth}{1.5pt}\color{primary}]
\titleformat{\subsection}{\large\bfseries\color{heading}}{}{0em}{}
\titleformat{\subsubsection}{\normalsize\bfseries\color{heading}}{}{0em}{}

% ─────────────────────────────────────────────────────────────────────────────
\begin{document}

% ── Title block ───────────────────────────────────────────────────────────────
\begin{center}
    \vspace*{0.8cm}
    {\Huge\bfseries\color{heading} Bachelor Thesis Expos\'{e}\par}
    \vspace{0.6cm}
    {\Large\bfseries\color{primary}
        Autonomous Mobile Manipulation:\\[0.2em]
        Design and Evaluation of a Pick-and-Place\\[0.2em]
        and Object Stacking Pipeline on the LIMO Cobot\par}
    \vspace{1.0cm}
    {\large Omar Alobaid\par}
    \vspace{0.2cm}
    {\large Supervisor: Prof.\ Ronny Hartanto\par}
    \vspace{0.2cm}
    {\large Rhine-Waal University of Applied Sciences (HSRW)\par}
    \vspace*{1.2cm}
\end{center}

% ─────────────────────────────────────────────────────────────────────────────
\section{Title}

\begin{tcolorbox}[colback=boxbg, colframe=primary, arc=3pt, boxrule=1.5pt,
    left=1em, right=1em, top=0.8em, bottom=0.8em]
    \textbf{Primary Title:}\\[0.4em]
    {\large Autonomous Mobile Manipulation: Design and Evaluation of a
    Pick-and-Place and Object Stacking Pipeline on the LIMO Cobot}
\end{tcolorbox}

\vspace{0.4em}
\noindent\textbf{Alternative titles:}
\begin{itemize}[label=\textbullet, leftmargin=1.5em, itemsep=0.1em]
    \item Autonomous Mobile Manipulation on the LIMO Cobot: Design, Evaluation, and
          Simulation-to-Reality Transfer of a Pick-and-Place Pipeline
    \item Autonomous Mobile Manipulation: Design and Evaluation of a Pick-and-Place
          Pipeline on a Table-Sized Mobile Manipulator (LIMO Cobot)
\end{itemize}

% ─────────────────────────────────────────────────────────────────────────────
\section{Research Question}

\begin{tcolorbox}[colback=boxbg, colframe=primary, arc=3pt, boxrule=1pt,
    left=1em, right=1em, top=1em, bottom=1em]
How can a modular autonomous mobile-manipulation pipeline, integrating onboard
navigation, perception, and collision-aware motion planning, be designed and evaluated
in a ROS\,2/Gazebo simulation to perform reliable, repeatable pick-and-place and
object-stacking tasks with a compact manipulator and gripper, and to what extent can
it then be transferred to and validated on the physical robot?
\end{tcolorbox}

\subsection{Sub-Questions}

\begin{enumerate}[label=\textbf{\arabic*.}, leftmargin=1.5em, itemsep=0.4em]

    \item \textbf{Navigation.}
    How reliably and repeatably can an autonomous mobile base position itself such that a
    target object lies within the manipulator's reachable workspace, and to what extent
    does base-positioning accuracy determine subsequent grasp success?

    \item \textbf{Perception.}
    How accurately and robustly can a target object's three-dimensional pose be estimated
    from onboard RGB-D sensing and expressed in the manipulator's planning frame, and how
    does the resulting estimation error propagate to grasp success?

    \item \textbf{Manipulation and reachability.}
    To what extent do the arm-mounting configuration, base stop distance, and workspace
    geometry constrain the reachable workspace of a short-reach manipulator, and how can
    collision-aware motion planning ensure that grasp and placement trajectories remain
    safe and executable within those constraints?

    \item \textbf{Stacking precision.}
    With what precision and repeatability can objects be placed to achieve stable
    sequential vertical stacking, and which sources of error most strongly limit stacking
    reliability?

    \item \textbf{Evaluation in simulation.}
    What level of task performance (success rate, placement accuracy, planning time, and
    cycle time) does the integrated pipeline achieve under repeated trials, and which
    failure modes dominate?

    \item \textbf{Simulation-to-reality transfer.}
    To what extent does a pipeline developed in simulation transfer to the physical robot,
    how does real-world performance compare with simulated performance, and what
    engineering effort does the transfer entail?

\end{enumerate}

\newpage

% ─────────────────────────────────────────────────────────────────────────────
\section{Scope and Limitations}

\subsection{Scope (Delimitations)}
\begin{itemize}[label=\checkmark, leftmargin=1.5em, itemsep=0.3em]
    \item The thesis designs and evaluates a complete mobile-manipulation pipeline
          covering navigation, perception, grasping, placement, and sequential stacking,
          for a single compact mobile manipulator (AgileX LIMO Cobot).
    \item The system is developed and primarily evaluated in physics-based simulation;
          deployment on the physical robot is pursued as a transfer-and-validation
          objective, contingent on hardware availability and the time available within
          the thesis period.
    \item The operating environment is \emph{static, planar, and structured}: a known
          layout, known support surfaces, and a single known target object at a time.
    \item Perception is delimited to a single, geometrically simple, colour-distinct
          object resting on a known support surface.
    \item Manipulation is restricted to top-down grasps of light, regularly shaped
          objects lying within the arm's reachable workspace.
    \item Evaluation is quantitative, based on success rate, placement accuracy, planning
          time, and cycle time, assessed under repeated trials ($N \geq 20$).
\end{itemize}

\subsection{Limitations}
\begin{itemize}[label=\textbullet, leftmargin=1.5em, itemsep=0.3em]
    \item Results are established primarily in simulation; despite physics-based
          modelling, a simulation-to-reality gap may affect real-world performance, and
          characterising this gap is itself part of the study.
    \item The limited reach of the manipulator ($\approx 0.28$\,m) constrains the
          feasible workspace and the admissible placement of objects and support surfaces.
    \item Classical colour-and-depth perception is sensitive to lighting, sensor noise,
          and object appearance; robustness under uncontrolled conditions is not claimed.
    \item Achievable stacking precision is bounded by the combined error of base
          positioning, perception, and end-effector alignment.
    \item Physical-hardware validation depends on platform availability, calibration, and
          the time remaining within the thesis period.
\end{itemize}

\subsection{Out of Scope}
\begin{tcolorbox}[colback=boxbg, colframe=accent, arc=3pt, boxrule=1pt,
    left=1em, right=1em, top=0.5em, bottom=0.5em]
    \begin{itemize}[label=\textbullet, leftmargin=1em, topsep=0.2em, itemsep=0.2em]
        \item Dynamic or deformable environments and moving obstacles/targets.
        \item Cluttered or multi-object scenes and general object recognition.
        \item Learning-based perception and learned manipulation policies.
        \item Multi-robot coordination.
        \item Online map-building (SLAM) during execution; a known map and layout are
              assumed.
    \end{itemize}
\end{tcolorbox}

\newpage

% ─────────────────────────────────────────────────────────────────────────────
\section{Task Planning}

\textbf{Total thesis period: 3 months (13 weeks).}
Phases overlap deliberately to make the most of the available time.
Continuous light drafting runs throughout Phases 1–4 so that the writing phase
is consolidation rather than writing from scratch.

\vspace{0.5em}

\noindent\textbf{Phase 1 — Consolidation (Weeks 1–2)}\\
Finalise the thesis framing; complete the integrated
navigate $\to$ dock $\to$ perceive $\to$ grasp $\to$ transport $\to$ place $\to$ stack
cycle into a reliable, repeatable run. Freeze the simulation testbed used for all
subsequent experiments. Begin literature review.

\vspace{0.5em}

\noindent\textbf{Phase 2 — Simulation Experiments (Weeks 2–6)}\\
Execute the systematic trial campaign ($N \geq 20$ per configuration), recording task
success rate, placement accuracy, planning time, and cycle time.
Conduct the reachability study (effect of mounting, base stop distance, and workspace
layout) and the stacking-precision study. Complete literature review (Chapter\,2).

\vspace{0.5em}

\noindent\textbf{Phase 3 — Simulation-to-Reality Transfer (Weeks 5–10, contingent)}\\
Set up the physical LIMO Cobot (onboard computer, drivers, arm and camera calibration),
port the pipeline to hardware, and perform the real-robot trials.
Compare real-world placement precision against simulation and record the engineering
effort required for the transfer.
\textit{Contingent on hardware availability and lab access.}

\vspace{0.5em}

\noindent\textbf{Phase 4 — Analysis (Weeks 9–11)}\\
Process and interpret collected data, produce figures and tables, identify dominant
failure modes, and prepare results for write-up.

\vspace{0.5em}

\noindent\textbf{Phase 5 — Dedicated Writing \& Submission (Weeks 10–13)}\\
A three-week period reserved for writing and final revision.
Complete all chapters, the discussion and conclusion, and the abstract;
finalise and submit.

\newpage

% ─────────────────────────────────────────────────────────────────────────────
\section{Gantt Chart}

\textbf{Total thesis period: 3 months (13 weeks)}

\vspace{1em}

\begin{center}
\begin{ganttchart}[
    hgrid,
    vgrid={*6{draw=none}, dotted},
    x unit=0.72cm,
    y unit title=0.6cm,
    y unit chart=0.65cm,
    title label font=\scriptsize,
    bar label font=\small,
    group label font=\small\bfseries,
    milestone label font=\small\itshape,
    bar/.append style={fill=ganttblue!60, rounded corners=1pt},
    bar incomplete/.append style={fill=ganttgrey!30},
    milestone/.append style={fill=accent, shape=diamond, inner sep=1.5pt},
    group/.append style={fill=ganttblue!20, draw=ganttblue},
    title/.append style={fill=boxbg},
    canvas/.append style={draw=border},
    today=1,
    today rule/.style={draw=accent, thick, dashed},
    today label=\scriptsize Today
]{1}{13}

    % ── Column headers (week numbers / months) ────────────────────────────────
    \gantttitle{Month 1}{4}
    \gantttitle{Month 2}{4}
    \gantttitle{Month 3}{4}
    \gantttitle{}{1} \\
    \gantttitlelist{1,...,13}{1} \\

    % ── Phase 1 ─────────────────────────────────────────────────────────────
    \ganttgroup{Phase 1 — Consolidation}{1}{2} \\
    \ganttbar{Freeze simulation testbed}{1}{2} \\
    \ganttbar{Literature review (begin)}{1}{2} \\

    % ── Phase 2 ─────────────────────────────────────────────────────────────
    \ganttgroup{Phase 2 — Simulation Experiments}{2}{6} \\
    \ganttbar{N\,$\geq$\,20 trial campaign}{2}{5} \\
    \ganttbar{Reachability \& stacking study}{4}{6} \\
    \ganttbar{Literature review (complete)}{3}{6} \\
    \ganttmilestone{M1 — Sim data collected}{6} \\

    % ── Phase 3 ─────────────────────────────────────────────────────────────
    \ganttgroup{Phase 3 — Sim-to-Real (contingent)}{5}{10} \\
    \ganttbar{Hardware setup \& calibration}{5}{7} \\
    \ganttbar{Real-robot trials}{7}{10} \\
    \ganttmilestone{M2 — Hardware trials done}{10} \\

    % ── Phase 4 ─────────────────────────────────────────────────────────────
    \ganttgroup{Phase 4 — Analysis}{9}{11} \\
    \ganttbar{Data processing \& figures}{9}{11} \\
    \ganttmilestone{M3 — Analysis complete}{11} \\

    % ── Phase 5 ─────────────────────────────────────────────────────────────
    \ganttgroup{Phase 5 — Writing \& Submission}{10}{13} \\
    \ganttbar{Chapters \& discussion}{10}{12} \\
    \ganttbar{Revision \& final checks}{12}{13} \\
    \ganttmilestone{M4 — Thesis submitted}{13}

\end{ganttchart}
\end{center}

\vspace{0.8em}

\subsection{Milestones}

\begin{table}[h]
\centering
\small
\begin{tabularx}{\textwidth}{l c X}
\toprule
\textbf{Milestone} & \textbf{Timing} & \textbf{Deliverable} \\
\midrule
M1 — Sim data collected   & End of Week 6  & $N\geq 20$ simulation runs complete; metrics logged. \\
M2 — Hardware trials done & End of Week 10 & Real-robot stacking trials complete; comparison data logged. \\
M3 — Analysis complete    & End of Week 11 & All figures, tables, and failure-mode analysis ready. \\
M4 — Thesis submitted     & End of Week 13 & Final corrected thesis submitted. \\
\bottomrule
\end{tabularx}
\end{table}

\newpage

% ─────────────────────────────────────────────────────────────────────────────
\section{Overview}

Mobile manipulation brings together two abilities that are normally studied apart:
moving through an environment and handling the objects in it.
A mobile manipulator is a wheeled base with a robotic arm, so one robot can drive to
an object, pick it up, carry it, and put it down.
The idea is well established and has been demonstrated on many research platforms over
the years~\cite{masannek2026eaiws}.
Making it work reliably on a small, low-cost robot is still demanding, because the base
and the arm have to be coordinated precisely for each task to succeed.

This thesis designs a complete pipeline for such a robot and studies how well it picks
objects up, places them, and stacks them.
The robot drives to a target, finds it with an onboard camera~\cite{chaumette2006visual},
plans a safe grasp around the surfaces near it, picks the object up, carries it, and
places it where it should go.
By repeating this cycle and setting each object on top of the last, the robot builds a
stack.
Stacking is the most demanding goal of the three, because the objects have to be placed
precisely enough to rest stably on one another, so it sets the standard the rest of the
pipeline has to meet.
The work is done first in simulation and then, as far as time and hardware allow, on
the real robot.
The methods are kept simple and transparent rather than learning-based, so the system
is easy to follow and to transfer to hardware.
Performance is judged by how often the task succeeds, how accurately objects are
placed, and how long planning and each full cycle take.
Comparing the simulated behaviour with the behaviour on the real robot is itself one of
the things the thesis sets out to measure.

\subsection{Platform Architecture}

The mobile manipulator platform combines an \textbf{AgileX LIMO} differential-drive
base with an \textbf{Elephant Robotics myCobot 280} 6-DoF arm (the LIMO Cobot
configuration).
The base provides autonomous navigation capabilities, while the arm extends
approximately 280\,mm with a payload capacity of 250\,g.
An adaptive parallel-jaw gripper is mounted at the end-effector, and an RGB-D camera
provides visual perception.
The entire system operates under \textbf{ROS\,2 Humble} with \textbf{Gazebo Classic}
as the simulation environment, enabling physics-based testing before hardware
deployment.
This compact configuration represents a low-cost platform for mobile manipulation
research, with the arm's limited reach creating stringent requirements for base
positioning accuracy.

\section{References}

\begin{itemize}[label=\textbullet, leftmargin=1.5em]
    \item Macenski, S.\ and Jambrecic, I., ``SLAM Toolbox: SLAM for the Dynamic World,''
          \textit{Journal of Open Source Software}, 2021.
    \item Macenski, S.\ et al., ``The Marathon 2: A Navigation System,'' \textit{IROS}, 2020.
    \item Masannek, M.\ et al., ``Embodied Autonomous Intelligence for On-Demand
          Manufacturing,'' EAI-WS Proposal, 2026.
    \item Chaumette, F.\ and Hutchinson, S., ``Visual Servo Control, Part I: Basic
          Approaches,'' \textit{IEEE Robotics \& Automation Magazine}, 2006.
\end{itemize}

\end{document}
